% 1
\frame[plain]{\titlepage}

% 2
\begin{frame}{Что мы уже знаем про градиентный спуск}
\small
Для $L$-гладкой функции при шаге $\alpha=1/L$:
\[
 x_{k+1}=x_k-\frac1L\grad f(x_k)
\]
и
\[
\boxed{
 f(x_{k+1})
 \le
 f(x_k)-\frac{1}{2L}\norm{\grad f(x_k)}^2.
}
\]
Следовательно, значения функции не возрастают.

\begin{alertblock}{Но этого мало}
Монотонное убывание ещё не говорит, \textbf{насколько быстро} мы приближаемся к оптимуму.
\end{alertblock}
\end{frame}

% 3
\begin{frame}{Из лекции 5: теперь мы знаем, куда хотим прийти}
\small
Для дифференцируемой выпуклой функции
\[
\grad f(x^\star)=0
\quad\Longrightarrow\quad
x^\star\text{ — глобальный минимум.}
\]
Для сильно выпуклой функции этот минимум единственен.

\vspace{0.8em}
Сегодня вопрос уже не в существовании хорошего решения, а в сложности его поиска:
\[
\boxed{\text{Сколько итераций нужно градиентному спуску?}}
\]
\end{frame}

% 4
\begin{frame}{Три разные меры ошибки}
\small
На итерации $k$ можно измерять качество по-разному:
\[
\boxed{f(x_k)-f^\star}
\qquad
\boxed{\norm{x_k-x^\star}}
\qquad
\boxed{\norm{\grad f(x_k)}}.
\]

\centering
\includegraphics[width=0.80\linewidth]{error_metrics.pdf}
\end{frame}

% 5
\begin{frame}{Они связаны, но не совпадают}
\small
Одинаковое расстояние до $x^\star$ может давать совершенно разную ошибку по функции — всё зависит от кривизны.

А маленький градиент сам по себе без дополнительных предположений не обязан означать маленькую ошибку по функции.

\begin{block}{Основная мера сегодня}
Будем оценивать
\[
\boxed{f(x_k)-f^\star.}
\]
Это непосредственно отвечает на вопрос, насколько хорошее значение целевой функции уже найдено.
\end{block}
\end{frame}

% 6
\begin{frame}{Два режима, которые мы получим}
\centering
\includegraphics[width=0.86\linewidth]{rate_preview.pdf}

\vspace{0.2em}
\small
Для выпуклой гладкой функции появится скорость порядка $1/k$, а сильная выпуклость даст геометрическое уменьшение ошибки.
\end{frame}

% 7
\begin{frame}{Главные результаты лекции — пока без доказательства}
\small
Для выпуклой $L$-гладкой функции:
\[
\boxed{
 f(x_k)-f^\star
 \le
 \frac{L\norm{x_0-x^\star}^2}{2k}.
}
\]
Для $\mu$-сильно выпуклой и $L$-гладкой:
\[
\boxed{
 f(x_k)-f^\star
 \le
 \left(1-\frac{\mu}{L}\right)^k
 \bigl(f(x_0)-f^\star\bigr).
}
\]
Наша задача — понять, откуда возникают обе формулы.
\end{frame}

% 8
\begin{frame}{$L$-гладкость: ограничение на максимальную кривизну}
\small
Напомним определение:
\[
\norm{\grad f(x)-\grad f(y)}
\le
L\norm{x-y}.
\]
Если $f$ дважды дифференцируема, то интуитивно
\[
L\approx\text{максимальная кривизна функции}.
\]
Для квадратичной функции
\[
f(x)=\frac12x^\top Ax-b^\top x
\]
имеем
\[
L=\lambda_{\max}(A).
\]
\end{frame}

% 9
\begin{frame}{Главный инструмент из лекции 4}
\small
Из $L$-гладкости следует квадратичная верхняя оценка:
\[
f(y)
\le
f(x)+\grad f(x)^\top(y-x)
+\frac L2\norm{y-x}^2.
\]
При градиентном шаге $y=x-\frac1L\grad f(x)$:
\[
\boxed{
 f(x_{k+1})
 \le
 f(x_k)-\frac1{2L}\norm{\grad f(x_k)}^2.
}
\]
\end{frame}

% 10
\begin{frame}{Главный инструмент из лекции 5}
\small
Для выпуклой дифференцируемой функции
\[
 f(y)
 \ge
 f(x)+\grad f(x)^\top(y-x).
\]
Это означает: касательная в $x$ — глобальная нижняя оценка функции.

\begin{block}{Нам понадобится одна специальная подстановка}
Положим $y=x^\star$.
\end{block}
\end{frame}

% 11
\begin{frame}{Выпуклость связывает градиент с ошибкой по функции}
\small
Подставляя $y=x^\star$, получаем
\[
f^\star
\ge
f(x)+\grad f(x)^\top(x^\star-x).
\]
После переноса:
\[
\boxed{
 f(x)-f^\star
 \le
 \grad f(x)^\top(x-x^\star).
}
\]
Справа появилась информация о направлении на оптимум.
\end{frame}

% 12
\begin{frame}{Сильная выпуклость добавляет квадратичный запас}
\small
Если $f$ $\mu$-сильно выпукла, то
\[
f(y)
\ge
f(x)+\grad f(x)^\top(y-x)
+\frac\mu2\norm{y-x}^2.
\]
Это более сильное утверждение, чем обычная выпуклость.

\centering
\includegraphics[width=0.62\linewidth]{strong_lower_model.pdf}
\end{frame}

% 13
\begin{frame}{Карта доказательств}
\small
В первой части:
\[
\boxed{L\text{-гладкость}+\text{выпуклость}}
\Longrightarrow
\boxed{f(x_k)-f^\star=O(1/k)}.
\]
Во второй части добавим сильную выпуклость:
\[
\boxed{L\text{-гладкость}+\mu\text{-сильная выпуклость}}
\Longrightarrow
\boxed{f(x_k)-f^\star\le q^k C}.
\]
И увидим, что
\[
q=1-\frac{1}{\kappa},
\qquad
\kappa=\frac{L}{\mu}.
\]
\end{frame}

% 14
\begin{frame}{Градиентный шаг даёт полезное равенство}
\small
Работаем с шагом
\[
x_{k+1}=x_k-\frac1L\grad f(x_k).
\]
Переносим:
\[
\frac1L\grad f(x_k)=x_k-x_{k+1}.
\]
Поэтому
\[
\boxed{
\grad f(x_k)=L(x_k-x_{k+1}).
}
\]
Эта простая формула позволит заменить градиент геометрией последовательных точек.
\end{frame}

% 15
\begin{frame}{Какой вид оценки нам хотелось бы получить?}
\small
Если удастся доказать неравенство вида
\[
\boxed{
 f(x_{k+1})-f^\star
 \le
 C\Bigl(
 \norm{x_k-x^\star}^2
 -
 \norm{x_{k+1}-x^\star}^2
 \Bigr),
}
\]
то при сложении по $k$ большая часть членов сократится.

\begin{alertblock}{Именно это и произойдёт}
Такие разности называются телескопическими: внутренние элементы исчезают при суммировании.
\end{alertblock}
\end{frame}
